%% ma: L10-B04-0004
%% lop: 10
%% chuong: 2
%% bai: 4
%% dang: 10.4.2
%% loai: TLN
%% muc: VD
%% dap_an: "34"
%% nguon: "(NBV)-ĐỀ SỐ 19-MỖI NGÀY 1 ĐỀ THI 2025.pdf (D19, MỖI NGÀY 1 ĐỀ - Đề số 19), Phần III Câu 1 (Chuyên Lương Thế Vinh – Đồng Nai 2025)"
%% kien_thuc: "Bài 4 (hệ bất phương trình bậc nhất hai ẩn, quy hoạch tuyến tính)"
%% trang_thai: chinh-thuc
%% ghi_chu: "Đáp án gốc 34 đúng (giải độc lập: các đỉnh (6;5), (9;0), (12;0), (12;10), (5;10); C(6;5) = 34 nhỏ nhất). Trùng ý tưởng với dạng đã duyệt 10.4.2 (quy hoạch tuyến tính tối thiểu chi phí với hệ bất phương trình bậc nhất hai ẩn), nhập chính thức. Bỏ ảnh minh họa xe chở lợn (ảnh trang trí)."
\begin{ex}
Một trang trại cần thuê xe vận chuyển $450$ con lợn và $35$ tấn cám. Nơi cho thuê xe chỉ có $12$ xe lớn và $10$ xe nhỏ. Một chiếc xe lớn có thể chở $50$ con lợn và $5$ tấn cám. Một chiếc xe nhỏ có thể chở $30$ con lợn và $1$ tấn cám. Tiền thuê một xe lớn là $4$ triệu đồng, một xe nhỏ là $2$ triệu đồng. Hỏi chi phí thuê xe ít nhất là bao nhiêu triệu đồng?
\shortans{34}
\loigiai{Gọi $x$, $y$ lần lượt là số xe lớn và số xe nhỏ được thuê ($x,y\in\mathbb N$). Ta có hệ
$$\begin{cases}0\le x\le12\\0\le y\le10\\50x+30y\ge450\\5x+y\ge35\end{cases}\Leftrightarrow\begin{cases}0\le x\le12\\0\le y\le10\\5x+3y\ge45\\5x+y\ge35.\end{cases}$$
Chi phí thuê xe là $C=4x+2y$ (triệu đồng). Miền nghiệm là ngũ giác có các đỉnh $(9;0)$, $(12;0)$, $(12;10)$, $(5;10)$, $(6;5)$, với $C(9;0)=36$, $C(12;0)=48$, $C(12;10)=68$, $C(5;10)=40$, $C(6;5)=34$. Vậy chi phí ít nhất là $34$ triệu đồng (thuê $6$ xe lớn và $5$ xe nhỏ).}
\end{ex}
